\documentclass[10pt]{article}
\usepackage[a4paper,margin=0.7in]{geometry}
\usepackage{ucu_doc}
\begin{document}
\doctitle{Safety-Jacket Follower Robot: User Manual}{Jamugisa Peter Paul (S25M25/008) and Ayebare Moses (S25M25/003) \textbar{} 30 October 2026}

\textbf{What it does.} A Hiwonder MasterPi robot finds the person wearing a hi-vis safety jacket, follows at about 1.5\,m, speeds up with them, searches when it loses them and stops when unsure. It runs in Gazebo simulation or on the real robot.

\section{1. What you need}
\begin{itemize}
  \item A Windows 11 laptop with an NVIDIA graphics card, and WSL~2 (Windows Subsystem for Linux 2) running Ubuntu 24.04. Without a graphics card, set \texttt{device} to \texttt{cpu} in \texttt{config/detector.yaml}.
  \item For the real robot: a MasterPi with \textbf{two fully charged 18650 cells} (about 8.2\,V) and 3\,m of clear floor.
\end{itemize}

\section{2. Install (once)}
Copy \texttt{6\_Source\_Code/jacket\_follower} to \texttt{D:\textbackslash mscs\textbackslash intelligent\_systems\_and\_robotics\textbackslash jacket\_follower} (the scripts expect this path; to use another, edit \texttt{PROJ} in the scripts). In PowerShell, run \texttt{wsl -d <your Ubuntu name> -u root}, then:
\begin{verbatim}
cd /mnt/d/mscs/intelligent_systems_and_robotics/jacket_follower/ros2_setup
bash install_ros2_jazzy.sh && bash install_ml_and_build.sh
\end{verbatim}
This installs ROS~2 (Robot Operating System 2) Jazzy, Gazebo Harmonic and the detector libraries, and builds the package.

\section{3. Run the simulation}
\begin{verbatim}
bash open_sim.sh s1_normal      # or s2_speed, or s3_challenge
bash open_sim.sh s3_challenge acquire_conf:=0.40   # failure case: lock-on rule off
\end{verbatim}
The Gazebo window and the camera view open. Watch from the start: people begin walking about 4\,s after loading, and the walk lasts about 1 minute. Press \textbf{Ctrl+C} to stop. Use the mouse only to look around; the move tools push the models. To repeat all measured runs: \texttt{bash run\_final.sh}, then \texttt{py -3 scripts/analyze\_sim.py} in Windows.

\section{4. Run the real robot}
\begin{enumerate}
  \item Switch on the battery box, then the expansion board. Wait about 1 minute, then join the robot's Wi-Fi hotspot on the laptop.
  \item First time only: from \texttt{robot/}, copy \texttt{masterpi\_server.py} and \texttt{start\_server.sh} to \texttt{/home/pi} on the robot. From \texttt{robot/tools/}, copy \texttt{follow\_pose.json}, \texttt{set\_follow\_pose\_slow.py} and \texttt{tilt\_up.py} there too.
  \item On the robot, run \texttt{bash /home/pi/start\_server.sh}, then set the arm with \texttt{sudo python3 set\_follow\_pose\_slow.py}.
  \item Test with the wheels off: \texttt{bash open\_real.sh}. A camera window shows the box and the state. Stand 1.5--2\,m away in the vest.
  \item When the box is steady, put the robot on the floor and run \texttt{bash open\_real.sh motion}.
\end{enumerate}

\section{5. Safety}
\begin{itemize}
  \item Keep people, cables and stairs clear of the robot. Stay within reach, and \textbf{lift the robot to stop it at any time}. Ctrl+C on the laptop also stops it.
  \item The robot stops itself within 0.5\,s if commands stop, and zeroes its motors at every boot (\texttt{motors\_off\_at\_boot.service}).
  \item Do not run below 7\,V: the Raspberry Pi can restart mid-run. Charge cells in the Hiwonder charger, not through the Pi's USB-C port.
\end{itemize}

\section{6. Troubleshooting}
\begin{tabularx}{\linewidth}{>{\raggedright\arraybackslash}p{5.2cm} X}
\toprule
\textbf{Problem} & \textbf{What to do} \\
\midrule
``Permission denied'' in WSL & You are not root. Type \texttt{exit}, then reopen with \texttt{wsl -d <name> -u root} \\
State stays E\_STOP (emergency stop), no picture & Laptop left the robot's Wi-Fi, or the server is not running: reconnect, rerun \texttt{start\_server.sh} \\
Robot keeps searching with the vest in view & Move closer (under 2.5\,m), add light, check focus; the vest may be too small or dim \\
Wheels do not move & Battery off or flat: check the battery-box switch and voltage \\
Vest at the top edge of the picture & Tilt the camera up: \texttt{sudo python3 tilt\_up.py 680} (a lower number tilts higher) \\
Wi-Fi vanishes during a run & Brown-out: switch off, fit charged cells \\
\bottomrule
\end{tabularx}
\end{document}
